\documentclass{article}
\usepackage{mathtools} 
\usepackage{fontspec}
\usepackage[UTF8]{ctex}
\usepackage{amsthm}
\usepackage{mdframed}
\usepackage{xcolor}
\usepackage{amssymb}
\usepackage{amsmath}


% 定义新的带灰色背景的说明环境 zremark
\newmdtheoremenv[
  backgroundcolor=gray!10,
  % 边框与背景一致，边框线会消失
  linecolor=gray!10
]{zremark}{说明}

% 通用矩阵命令: \flexmatrix{矩阵名}{元素符号}{行数}{列数}
\newcommand{\flexmatrix}[4]{
  \[
  #1 = \begin{pmatrix}
    #2_{11}     & #2_{12}     & \cdots & #2_{1#4}   \\
    #2_{21}     & #2_{22}     & \cdots & #2_{2#4}   \\
    \vdots      & \vdots      & \ddots & \vdots     \\
    #2_{#31}    & #2_{#32}    & \cdots & #2_{#3#4}
  \end{pmatrix}
  \]
}

% 简化版命令（默认矩阵名为A，元素符号为a）: \quickmatrix{行数}{列数}
\newcommand{\quickmatrix}[2]{\flexmatrix{A}{a}{#1}{#2}}

\begin{document}
\title{习题 13.1}
\author{张志聪}
\maketitle

\section*{1}
\begin{itemize}
  \item (1)

        对任意$x \in (-1, 1)$，我们有
        \begin{align*}
          \lim\limits_{n \to \infty} \sqrt{x^2 - \frac{1}{n^2}}
           & = \sqrt{x^2}
        \end{align*}
        令$f(x) = \sqrt{x^2}$，于是$f(x)$是函数列的极限函数。

        我们有
        \begin{align*}
          \lim\limits_{n \to \infty} \sup\limits_{x \in (-1, 1)} |f_n(x) - f(x)|
           & = \lim\limits_{n \to \infty} \sup\limits_{x \in (-1, 1)} \left|\sqrt{x^2 - \frac{1}{n^2}} - \sqrt{x^2}\right|                                   \\
           & = \lim\limits_{n \to \infty} \sup\limits_{x \in (-1, 1)} \left|\frac{x^2 - \frac{1}{n^2} - x^2}{\sqrt{x^2 - \frac{1}{n^2}} + \sqrt{x^2}}\right| \\
           & = \lim\limits_{n \to \infty} \sup\limits_{x \in (-1, 1)} \left|\frac{\frac{1}{n^2}}{\sqrt{x^2 - \frac{1}{n^2}} + \sqrt{x^2}}\right|
        \end{align*}
        可知$x = 0$时，取最大值，于是
        \begin{align*}
          \lim\limits_{n \to \infty} \sup\limits_{x \in (-1, 1)} \left|\frac{\frac{1}{n^2}}{\sqrt{x^2 - \frac{1}{n^2}} + \sqrt{x^2}}\right|
           & = \lim\limits_{n \to \infty} \sup\limits_{x \in (-1, 1)} \frac{1}{n} \\
           & = 0
        \end{align*}
        故$f_n(x) \rightrightarrows f(x)$.

  \item (2)

        对任意$x \in (-\infty, +\infty)$，我们有
        \begin{align*}
          \lim\limits_{n \to \infty} \frac{x}{1 + n^2x^2}
           & = 0
        \end{align*}
        令$f(x) = 0$，于是$f(x)$是函数列的极限函数。

        考虑
        \begin{align*}
          |f_n(x) - f(x)|
           & = \left|\frac{x}{1 + n^2x^2}\right|                 \\
           & = \frac{|x|}{1 + n^2x^2}                            \\
           & = \frac{1}{\frac{1}{|x|} + n^2|x|}                  \\
           & \leq \frac{1}{ 2 \sqrt{\frac{1}{|x|} \cdot n^2|x|}} \\
           & = \frac{1}{ 2 n}
        \end{align*}
        于是，我们有
        \begin{align*}
          0 \leq \sup\limits_{x \in (-\infty, +\infty)} |f_n(x) - f(x)| \leq \frac{1}{2n}
        \end{align*}
        由迫敛性可知
        \begin{align*}
          \lim\limits_{n \to \infty} \sup\limits_{x \in (-\infty, +\infty)} |f_n(x) - f(x)| = 0
        \end{align*}

  \item (3)

        $x = 0$时，$f_n (x) = 1, n = 1, 2, 3, \cdots$；

        我们有
        \begin{align*}
          \lim\limits_{n \to \infty} \frac{1}{n + 1} = 0
        \end{align*}
        于是，对任意$x \in (0, 1)$，存在$N$，当$n > N$，有
        \begin{align*}
          \frac{1}{n + 1} & < x \\
          f_n(x)          & = 0
        \end{align*}
        所以
        \begin{align*}
          \lim\limits_{n \to \infty} f_n(x)
           & = 0
        \end{align*}
        令
        \begin{equation*}
          f(x) =
          \begin{cases}
            1, & x = 0        \\
            0, & x \in (0, 1)
          \end{cases}
        \end{equation*}
        $f(x)$是函数列的极限函数。

        \begin{itemize}
          \item 是否为一致收敛；

                取$x = \frac{1}{2(n + 1)}$，有
                \begin{align*}
                  |f_n(x) - f(x)|
                   & = \left|-(n+1) \frac{1}{2(n + 1)} + 1\right| \\
                   & = \left|-\frac{1}{2} + 1\right|              \\
                   & = \frac{1}{2}
                \end{align*}
                由上确界的定义可知，
                \begin{align*}
                  \sup\limits_{x \in [0, 1)} |f_n(x) - f(x)| \geq \frac{1}{2}
                \end{align*}
                所以由保不等式性
                \begin{align*}
                  \lim\limits_{n \to \infty} \sup\limits_{x \in [0, 1)} |f_n(x) - f(x)| \neq 0
                \end{align*}
                故不一致收敛。

          \item 是否内闭一致收敛。

                取闭区间为$[0, \frac{1}{2}]$，证明同上，即不内闭一致收敛。
        \end{itemize}

        \begin{zremark}
          扩充：

          区间$[0, 1)$改为$(0, 1)$后，会有什么变化。
        \end{zremark}

        \begin{itemize}
          \item 一致收敛性；

                同理，可取$x = \frac{1}{2(n + 1)}$，来证明一致收敛性。

                原因在于：影响一致收敛的是点$0$的领域，而不是这一个点。

          \item 内闭一致收敛性。

                设$[a, b] \subset (0, 1)$，由$\lim\limits_{n \to \infty} \frac{1}{n + 1} = 0$，
                当$n$足够大时，由$\frac{1}{n + 1} < a$，此时$f_n(x)$在$[a, b]$上不在是分段函数，
                $f_n(x) \equiv 0$，显然，内闭一致收敛。

                原因在于：闭区间和开区间有着巨大的差异，闭区间$[a, b]$不要按照去掉一个点$0$理解，
                而是去掉了$U(0; a)$这个邻域。
        \end{itemize}

  \item (4)

        对任意$x \in D$，我们有
        \begin{align*}
          \lim\limits_{n \to \infty} \frac{x}{n} = 0
        \end{align*}
        令$f(x) = 0$，于是$f(x)$是函数列的极限函数。

        \begin{itemize}
          \item 是否为一致收敛；

                考虑
                \begin{align*}
                  |f_n(x) - f(x)|
                   & = f_n(x)      \\
                   & = \frac{x}{n}
                \end{align*}
                因为
                \begin{align*}
                  \lim\limits_{x \to \infty} \frac{x}{n} = +\infty
                \end{align*}
                于是对$M > 1$，存在$x_0$，$x > x_0$，有
                \begin{align*}
                  \frac{x}{n} > M
                \end{align*}
                可得
                \begin{align*}
                  \sup \limits_{x \in D} |f_n(x) - f(x)| > M
                \end{align*}
                从而
                \begin{align*}
                  \lim\limits_{n \to \infty} \sup\limits_{x \in D} |f_n(x) - f(x)| \neq 0
                \end{align*}
                故不一致收敛。

          \item 是否内闭一致收敛。

                对任意$[a, b] \subset D$，$f_n(x)$在区间$[a, b]$上单增，可知
                \begin{align*}
                  \sup\limits_{x \in [a, b]} |f_n(x) - f(x)| = \frac{b}{n}
                \end{align*}
                于是
                \begin{align*}
                  \lim\limits_{n \to \infty} \sup\limits_{x \in [a, b]} |f_n(x) - f(x)|
                   & = \lim\limits_{n \to \infty} \frac{b}{n} \\
                   & = 0
                \end{align*}
                故内闭一致收敛。
        \end{itemize}

  \item (5)

        对任意$x \in D$，我们有
        \begin{align*}
          \lim\limits_{n \to \infty} sin \frac{x}{n} = 0
        \end{align*}
        令$f(x) = 0$，于是$f(x)$是函数列的极限函数。

        \begin{itemize}
          \item 是否为一致收敛；

                考虑
                \begin{align*}
                  |f_n(x) - f(x)|
                   & = |f_n(x)|                     \\
                   & = \left|sin \frac{x}{n}\right|
                \end{align*}
                取$x_0 = \frac{n\pi}{2}$，
                \begin{align*}
                  \left|sin \frac{x_0}{n}\right| = 1
                \end{align*}
                所以
                \begin{align*}
                  \sup\limits_{x \in D} |f_n(x) - f(x)| \geq 1
                \end{align*}
                于是，我们有
                \begin{align*}
                  \lim\limits_{n \to \infty} \sup\limits_{x \in D} |f_n(x) - f(x)| \neq 0
                \end{align*}
                故不一致收敛。

          \item 是否为内闭一致收敛。

                任意$[a, b] \in D$，
                我们有
                \begin{align*}
                  |f_n(x) - f(x)|
                   & = |f_n(x)|                     \\
                   & = \left|sin \frac{x}{n}\right| \\
                   & \leq \left| \frac{x}{n}\right|
                \end{align*}
                所以
                \begin{align*}
                  \sup\limits_{x \in [a, b]} |f_n(x) - f(x)| \leq max\{\frac{|a|}{n},\frac{|b|}{n}\}
                \end{align*}
                我们有
                \begin{align*}
                  \lim\limits_{n \to \infty} \sup\limits_{x \in [a, b]} |f_n(x) - f(x)|
                   & = \lim\limits_{n \to \infty} max\{\frac{|a|}{n},\frac{|b|}{n}\} \\
                   & = 0
                \end{align*}
                故内闭一致收敛。
        \end{itemize}
\end{itemize}

\section*{2}

对任意$n$有
\begin{align*}
  |f_n(x) - f(x)| \leq a_n
\end{align*}
即
\begin{align*}
  \sup\limits_{x \in D} |f_n(x) - f(x)| \leq a_n
\end{align*}
已知
\begin{align*}
  \lim\limits_{n \to \infty} a_n = 0
\end{align*}
由迫敛性可知
\begin{align*}
  \lim\limits_{n \to \infty} \sup\limits_{x \in D} |f_n(x) - f(x)| = 0
\end{align*}
故函数列一致收敛于$f$。

\section*{3}

\begin{itemize}
  \item (1)
        \begin{align*}
          |u_n|
           & = \left| \frac{x^n}{(n - 1)!} \right| \\
           & = \frac{|x|^n}{(n - 1)!}              \\
           & \leq \frac{r^n}{(n - 1)!}
        \end{align*}
        令$M_n = \frac{r^n}{(n - 1)!}$.
        \begin{align*}
          \lim\limits_{n \to \infty} \frac{M_{n + 1}}{M_{n}}
           & = \lim\limits_{n \to \infty} \frac{r^{n + 1}}{n!}  \frac{(n - 1)!}{r^n} \\
           & = \lim\limits_{n \to \infty} \frac{r}{n}                                \\
           & = 0
        \end{align*}
        由比式判别法可知
        \begin{align*}
          \sum M_n
        \end{align*}
        收敛。
        于是，由优级数判别法可知
        \begin{align*}
          \sum \frac{x^n}{(n - 1)!}
        \end{align*}
        收敛。

  \item (2)

        使用$A-D$判别法。

        取$\{v_n(x)\} = \{\frac{x^2}{(1 + x^2)^n} \}, \sum u_n(x) = \sum (-1)^{n - 1}$。

        (i) $\sum u_n(x)$的部分和函数列
        \begin{align*}
          |S_n(x)| \leq 1
        \end{align*}
        即$\sum u_n(x)$在$(-\infty, +\infty)$上一致有界。

        (ii) 对每一个$x \in (-\infty, +\infty)$，
        \begin{align*}
          (1 + x^2)^n \leq (1 + x^2)^{n + 1}
        \end{align*}
        所以
        \begin{align*}
          \frac{x^2}{(1 + x^2)^n} \geq \frac{1}{(1 + x^2)^{n + 1}}
        \end{align*}
        所以$\{v_n(x)\}$是单调的；

        (iii)对任意$x \in (-\infty, +\infty)$，
        \begin{align*}
          \lim\limits_{n \to \infty} \frac{x^2}{(1 + x^2)^n} = 0
        \end{align*}
        令$f(x) = 0$，于是$f(x)$是函数列的极限函数。

        考虑$x \neq 0$时，
        \begin{align*}
          |v_n(x) - f(x)|
           & = \left|\frac{x^2}{(1 + x^2)^n} - 0\right| \\
           & = \left|\frac{x^2}{(1 + x^2)^n}\right|     \\
           & \leq \frac{x^2}{1 + n x^2}                 \\
           & = \frac{1}{\frac{1}{x^2} + n}              \\
           & \leq \frac{1}{n}
        \end{align*}
        注：第一个不等式可以使用二次项展开得到。

        $x = 0$时，$|v_n(x) - f(x)| = 0 < \frac{1}{n}$；所以
        \begin{align*}
          \sup\limits_{x \in (-\infty, \infty)} |v_n(x) - f(x)| \leq \frac{1}{n}
        \end{align*}
        又有
        \begin{align*}
          \lim\limits_{n \to \infty} \sup\limits_{x \in (-\infty, \infty)} |v_n(x) - f(x)|
           & = \lim\limits_{n \to \infty} \frac{1}{n} \\
           & = 0
        \end{align*}
        所以，$v_n(x) \rightrightarrows 0$。

        综上，使用$A-D$判别法，函数项级数$\sum \frac{(-1)^{n - 1}x^2}{(1 + x^2)^n}$一致收敛。

\end{itemize}


\end{document}